\section{Defects and two-level systems}\label{app:E.22}

In fused silica a group of atoms can occupy either of two configurations, a fraction of an \r{a}ngstr\"om apart and equal in energy to within a meV, and the group absorbs and scatters at micro-electronvolt energies. A close Frenkel pair in silicon stores a few eV and may give them back long after the event that made it. \Cref{sec:16.7,sec:16.8,sec:19.6} describe both kinds of defect with a few numbers. For defects in crystals the numbers are the ohmic coupling $\alpha_{\mathrm{ohm}}$ to the electrons of a metal, the tunneling exponents and Huang--Rhys factors of release paths, and the barriers and masses of crowdions and dumbbells. For glasses the numbers are the density vertex of a two-level system and the distribution of two-level parameters. None of these numbers has been computed for a real material at the level those sections need. Each calculation below is built from the methods of \cref{app:E.10,app:E.11}: relaxed supercells, nudged elastic bands, density-functional perturbation theory and molecular dynamics on machine-learned potentials. \Cref{tab:E22-summary} collects the targets, the tests and the costs.

\subsection{Ohmic coupling of tunneling defects}

A defect that hops between two configurations L and R in a normal metal changes the potential seen by the conduction electrons, and the Fermi sea has to rearrange at every hop. \Cref{sec:16.8} turns the rearrangement of the Fermi sea into a relaxation rate through the ohmic coupling $\alpha_{\mathrm{ohm}}$. We compute the dimensionless $\alpha_{\mathrm{ohm}}$ of \cref{eq:E.10-ohm} from the phase-shift differences $\tilde\delta_\nu$, the eigenphases of $S_RS_L^\dagger$ at the Fermi energy. The targets are a crowdion in tungsten, hydrogen in niobium, and two-level systems in aluminum, its oxide and amorphous WSi. On the truncation label $\mathfrak{T}$ (\cref{sec:8.5}) the calculation is at the base point for the electrons, with Kohn--Sham scattering states at clamped nuclei, $n_B = 0$; the defect enters only through its two relaxed geometries.

Two routes give the phases. The first computes $S_L$ and $S_R$ directly, from the impurity Green's function of a Korringa--Kohn--Rostoker code (JuKKR or AKAIKKR, full potential, angular momenta up to $l = 3$ or 4) with each relaxed configuration embedded in the host. The two matrices are expressed in one channel basis centered at the midpoint of the hop, and \cref{eq:E.10-ohm} is applied to the eigenphases. A basis centered on the defect is not enough: a hop displaces the scatterer, so for crowdions and hydrogen the eigenphase form holds and a single-center phase-shift sum does not.

The second route uses plane-wave supercells and the orthogonality catastrophe itself. The many-electron ground states of the two configurations in a cell of size $L$ overlap as $|\langle\Psi_L|\Psi_R\rangle| \propto (k_{\mathrm{F}}L)^{-\frac12\sum_{\sigma\nu}(\tilde\delta_\nu/\pi)^2}$~\cite{anderson1967infrared}, and by \cref{eq:E.10-ohm} the exponent is $\tfrac12\sum_{\sigma\nu}\tilde\delta_\nu^2/\pi^2 = \alpha_{\mathrm{ohm}}$. The slope of $\ln|\langle\Psi_L|\Psi_R\rangle|$ against $\ln L$ is therefore $-\alpha_{\mathrm{ohm}}$. The overlap is the determinant of the overlap matrix between the occupied Kohn--Sham orbitals of the two cells. For tungsten we take cubic cells of 54, 128 and 250 atoms ($3a$, $4a$ and $5a$) with the PBE functional, an optimized norm-conserving potential with the $5s$ and $5p$ semicore states and a cutoff of \SI{60}{Ry}, all at $\mathrm{\Gamma}$ only. The length $L$ of the slope is then the cell side itself; a $\vect{k}$ grid that shrinks as the cell grows would hold the Born--von K\'arm\'an length fixed and leave the slope undefined. A sparser fixed-density grid, the same in every cell, is the alternative, with $L$ its Born--von K\'arm\'an length. The 250-atom cell holds about $2.1\times10^5$ plane waves and 2100 bands, about \SI{7}{GB} per $\vect{k}$ point. Each defect costs about $10^3$ hours on graphics processors.

A contact coupling $\pm V_{\mathrm{d}}$ in a single channel returns $\alpha_{\mathrm{ohm}} = 4K_{\mathrm{b}}$ (\cref{app:E.10}); the same code also reproduces the exponent of a one-dimensional tight-binding chain with one impurity site, whose phase shift is known in closed form. The measurement is the damping of aluminum-covered nanomechanical resonators~\cite{maillet2023nanomechanical}. The computed $\alpha_{\mathrm{ohm}}$ contains every order in the phase shifts, and it is compared through the rate of \cref{sec:16.8} with the measured damping, not with the Born reading of the fitted couplings. $\vect{k}$ sampling and smearing blur the Fermi surface on which the logarithm rests, so the overlaps must be taken with fixed occupations at a common Fermi level and the smearing extrapolated to zero. The host must be relaxed identically in both cells except near the defect, or the overlap measures the relaxation and not the hop. The computed couplings set the electron-assisted relaxation rates of \cref{sec:16.8} and, for tungsten, the coherence of crowdion motion in \cref{sec:19.6}.

\subsection{Release paths of eV-scale defects}

Two quantities decide whether a close Frenkel pair in silicon or germanium, or a split interstitial in aluminum, releases its stored energy at millikelvin: the tunneling exponent through its barrier, and the number of local vibrations that accept the eV released in the hop (\cref{sec:19.6,sec:16.8}). Let $\Delta\vect{R}_I$ be the change of the position of atom $I$ between the two quenched minima, and $\vect{e}_{k,I}$ the mass-weighted normal modes, of energies $\omega_k$, at the final minimum. The displacement of mode $k$ is $\Delta Q_k = \sum_I M_I^{1/2}\vect{e}_{k,I}\cdot\Delta\vect{R}_I$. A harmonic oscillator displaced by $\Delta Q_k$ stores $\tfrac12\omega_k^2\Delta Q_k^2$, which in units of its quantum is the Huang--Rhys factor~\cite{huang1950theory}
\begin{equation}
  S_{\mathrm{HR},k} = \tfrac12\,\omega_k\,\Delta Q_k^2 ,
  \qquad E_{\mathrm{reorg}} = \sum_k S_{\mathrm{HR},k}\,\omega_k .
  \label{eq:E.22-SHR}
\end{equation}
These quantities lie on the Born--Oppenheimer surface, $n_B = 0$, harmonic about each minimum, with the hop treated either as a two-state transition or, when the reduction fails, by an instanton on the full surface. The reorganization energy $E_{\mathrm{reorg}}$ decides the regime. If it is small compared with the energy $\varepsilon_{\mathrm{hop}}$ released in the hop, the hop must put that energy into many quanta of weakly coupled modes, and the rate is suppressed by the energy-gap law. If $E_{\mathrm{reorg}} \approx \varepsilon_{\mathrm{hop}}$, the reaction coordinate itself takes the energy and the energy-gap suppression is absent. The tunneling exponent is $\lambda_{\mathrm{tun}} = \int\dd s\,[2(V(s) - E)]^{1/2}$ along the mass-weighted minimum-energy path $V(s)$, and the zero-temperature ring-polymer instanton action equals $2\lambda_{\mathrm{tun}}$ when the path does not cut corners~\cite{lawrence2024semiclassical}.

A release within a time $t_{\mathrm{obs}}$ needs the tunneling factor $\ee^{-2\lambda_{\mathrm{tun}}}$ to beat $1/(\varGamma_0t_{\mathrm{obs}})$, with $\varGamma_0$ the attempt rate of \cref{app:E.11}. For a silicon close pair in one year the condition is $\lambda_{\mathrm{tun}} \le 23.6$ in the strong-coupling regime, against a two-parabola estimate of 27.5 (\cref{sec:19.6}), so $\lambda_{\mathrm{tun}}$ and $E_{\mathrm{reorg}}/\varepsilon_{\mathrm{hop}}$ of real paths are needed to about \SI{5}{\percent}.

A first calculation uses a 64-atom silicon cell with a $\langle110\rangle$ split interstitial 2 to \SI{4}{\angstrom} from a vacancy. It uses Quantum ESPRESSO with PBE and optimized norm-conserving potentials, a cutoff of \SI{40}{Ry} checked at 30 and \SI{50}{Ry}, a $2^3$ Monkhorst--Pack grid, and the cell fixed at the PBE lattice constant. The climbing-image nudged elastic band of \cref{app:E.11} runs with nine images to a force tolerance of \SI{5}{meV/\angstrom}. The Hessians at both minima come from finite displacements of \SI{0.01}{\angstrom}, $64\times3\times2 = 384$ calculations each. The path then gives $S_{\mathrm{HR},k}$, $E_{\mathrm{reorg}}$, $\lambda_{\mathrm{tun}}$ and the line shape of the hop, and one ring-polymer instanton with 64 to 128 beads, run with i-PI on the machine-learned potential, checks the action against $2\lambda_{\mathrm{tun}}$. The cost of this first calculation is set by the roughly 900 self-consistent calculations, about 120 core-hours, at under \SI{0.1}{GB} per $\vect{k}$ point. The full catalogue needs 50 paths per material, mapped on the machine-learned potential and checked against density-functional theory. A kinetic Monte Carlo over the resulting rates at 10 to \SI{100}{mK}, with $10^6$ seeds, gives the release spectrum. The cost of the catalogue, mostly the instantons, is estimated at $10^4$ to $10^5$ hours on graphics processors.

The barrier and the exponent are converged in the cutoff at 30, 40 and \SI{50}{Ry} and in the grid at $2^3$ and $3^3$, and one path is repeated in a 216-atom cell, since the 64-atom cell compresses the strain field of the pair. The catalogue is accepted by the release-path sum rule of \cref{sec:19.6}: the energy released by all hops of the kinetic Monte Carlo, plus the energy still stored at the end, equals the initial stored energy, with every energy taken from the table of configurational energies of \cref{app:E.23}. Charged close pairs need the finite-size corrections of \cref{app:E.10}. A charge-state variant of the same path, computed by constrained occupations, gives the Huang--Rhys factor $S_{\mathrm{HR},R}$ and frequency $\omega_R$ of the reaction coordinate when a carrier is captured. The probability that a capture drives the reaction over a barrier $E_{\mathrm{b}}$ is then the tail above $n = E_{\mathrm{b}}/\omega_R$ of the Poisson series of \cref{eq:E.15-ste}. Carrier capture is thus the athermal channel by which the carriers of a later event can release stored energy (\cref{sec:19.6}).

\subsection{Crowdions and dumbbells}

A crowdion in tungsten is an extra atom spread along a $\langle111\rangle$ row. It glides over a Peierls barrier of meV in atomistic calculations and of micro-electronvolts in the Frenkel--Kontorova model~\cite{fitzgerald2018structure}, and its kink mass $M_{\mathrm{cr}}$ is a fraction of an atomic mass (\cref{app:C.20}). \Cref{sec:19.6} asks whether the crowdion moves coherently at millikelvin. Coherent motion needs a band that survives three effects: the phonons that distinguish neighboring sites, the conduction electrons of the normal metal (bulk tungsten becomes superconducting only near \SI{15}{mK}), and the strain field of the vacancy it left behind. The phonon dressing multiplies the hopping by $\ee^{-S_{\mathrm{pol}}(T)}$, with
\begin{equation}
  S_{\mathrm{pol}}(T) = \frac{1}{\pi}\int\dd\omega\,\frac{J_{\mathrm{s}}(\omega)}{\omega^2}\coth\frac{\omega}{2T} ,
  \label{eq:E.22-Spol}
\end{equation}
for a linear coupling with spectral density $J_{\mathrm{s}}$ of \cref{app:C.20}. For acoustic phonons $J_{\mathrm{s}} = J_{\mathrm{s},3}\omega^3$ up to the Debye energy $\omega_{\mathrm{D}}$, which gives a finite $S_{\mathrm{pol}}(0) = J_{\mathrm{s},3}\omega_{\mathrm{D}}^2/2\pi$ and a thermal part $\pi J_{\mathrm{s},3}T^2/3$ (from $\int_0^\infty\dd x\,x/(\ee^x - 1) = \pi^2/6$). The electrons renormalize the hopping further by a power of $\alpha_{\mathrm{ohm}}$ (\cref{app:C.20}), with $\alpha_{\mathrm{ohm}}$ computed as above for the crowdion geometry.

The crowdion is treated on the Born--Oppenheimer surface, $n_B = 0$, with phonons and electrons coupled linearly to its position. The calculation for tungsten uses PBE, the semicore potential above, a 250-atom body-centered cubic cell, a cutoff of \SI{50}{Ry}, a $3^3$ grid and Marzari--Vanderbilt smearing of \SI{0.01}{Ry}. A climbing-image nudged elastic band along $\langle111\rangle$ gives the Peierls profile $E(z_0)$ of the kink center. The Hessian of the crowdion core gives the phonon coupling and $J_{\mathrm{s}}$. A fit of the profile to the Frenkel--Kontorova form gives the kink width and mass. The calculation also maps the interaction energy between the crowdion and its parent vacancy as a function of their separation. This map gives the radius inside which the strain gradient across one spacing exceeds the band width and the crowdion is bound. For aluminum the same steps give the migration of the $\langle100\rangle$ dumbbell, in cells of 108 to 256 atoms, and the one-dimensional glide of di- and tri-interstitials comes from machine-learned molecular dynamics at 10 to \SI{100}{K}, extrapolated to low temperature with instanton rates. Each material costs about $10^5$ core-hours.

The Peierls barrier is below a meV per cell, so the energies are differences at identical meshes and smearing, extrapolated to zero smearing and checked on a $4^3$ grid, with noise below \SI{0.1}{meV} per cell. The fit to the Frenkel--Kontorova form is valid only if the profile is sinusoidal; its higher harmonics measure the failure, and the kink width over the row spacing controls the continuum description (\cref{app:C.20}). In \cref{sec:19.6} the barrier and mass of the kink replace the model values, and the polaron factor, the electron renormalization and the confinement radius decide whether the crowdion band is real.

\subsection{Density vertex of a two-level system}

Dark matter couples to a two-level system in a glass through the change $f_{\mathrm{TLS}}(\vect{q})$ of the electron density between its configurations, \cref{eq:16.7-f}, which enters the scattering rate of \cref{sec:16.7} as $|f_{\mathrm{TLS}}|^2$. The estimate there moves one oxygen ion, eight electrons, rigidly by \SI{1}{\angstrom}, and misses the redistribution of the valence electrons, the screening of the dipole by the glass and the variety of real geometries. The vertex is a ground-state density at $n_B = 0$, with the adiabatic electrons following the atoms; the screening is that of the self-consistent density, and no further dielectric factor is applied.

The preferred method takes two self-consistent calculations of the same glass cell, relaxed in configurations L and R, and Fourier transforms the difference of the two densities at the reciprocal vectors $\vect{G}_{\mathrm{sc}}$ of the cell. It is exact to first order in the dark-matter coupling and keeps every order in the atomic displacement $\vect{d}$. Keeping every order matters because $qd$ reaches 0.5 to 1 at $q = 1$ to \SI{2}{keV} for $d = \SI{1}{\angstrom}$, where a linearization errs at order $(qd)^2$. The alternative is density-functional perturbation theory at $\Gamma$ of the supercell, with only the core atoms displaced. Every $\vect{G}_{\mathrm{sc}}$ is equivalent to $\Gamma$, so one response gives the density change at all of them, and the vertex is the sum over core atoms of the density response to each atom times its displacement from L to R. This route costs about ten times more and serves as a check on the relaxation field of the host.

The cells hold 216 atoms (72 formula units) of amorphous SiO$_2$, a cube of side \SI{15}{\angstrom}. Each two-level system comes from the large glasses of the ensemble below and is carved out with its environment to that size; the outer shell is re-quenched with the core held fixed, and the two configurations are re-relaxed by a nudged elastic band in the small cell; systems whose asymmetry moves by more than a meV in the transplant are discarded. The calculation uses PBE, optimized norm-conserving potentials, a wavefunction cutoff of \SI{80}{Ry} checked at 70 and \SI{90}{Ry}, a density cutoff of \SI{320}{Ry}, a self-consistency threshold of $10^{-10}$ and the $\Gamma$ point only. The momentum transfers run from $2\pi/L = \SI{0.42}{\angstrom^{-1}}$ (\SI{0.83}{keV}) to \SI{5}{\angstrom^{-1}}, and the 576 occupied bands take about \SI{2.5}{GB}. Each self-consistent calculation costs about one hour on a graphics processor, so 100 two-level systems cost about 200 hours with the preferred method.

The vertex is converged when it changes by less than \SI{2}{\percent} with the cutoff raised by \SI{10}{Ry} and with the cell enlarged from 216 to 648 atoms for one system; the ensemble average $\langle|f_{\mathrm{TLS}}|^2\rangle$ needs at least 100 systems for \SI{10}{\percent}. The benchmark is the rigid-ion value, which the computed vertex approaches at large $q$, where only the core electrons of the moving atoms contribute. Below $q = 2\pi/L$ the periodic cell forces the macroscopic field to vanish, so the cell gives the unscreened dipole of the core. The long-range part must be added analytically at small $q$; it is the change of the electronic dipole, screened by the high-frequency dielectric constant of the glass. Rebuilding the density change from Born effective charges places point charges where the vertex needs the valence density. \Cref{sec:16.7} uses the ratio of the computed $\langle|f_{\mathrm{TLS}}|^2\rangle$ to the rigid-ion value, together with the interference between core and host for cores off the center of the defect, which vanishes only for symmetric cores.

\subsection{Ensemble sampling of two-level systems}

The standard tunneling model of \cref{sec:16.7} assumes a flat distribution of asymmetries and of tunneling exponents, with density $\bar P$, and \cref{sec:16.8} builds its relaxation rates on this distribution. Testing that assumption for a specific glass needs the distribution itself, computed from the potential-energy surface of the glass. The targets are amorphous SiO$_2$, amorphous Al$_2$O$_3$ and amorphous WSi. For each we need the distributions of the asymmetry $\Delta_{\mathrm{as}}$, the tunneling amplitude $\Delta_0$ of \cref{app:C.20}, the deformation potential and the dipole moment. From them follow tables of the two-level response of \cref{sec:16.7} on at least 590 directions of $\vect{q}$, 60 logarithmic points in $|\vect{q}|$ from \SI{0.1}{eV} to \SI{10}{keV}, and 50 energies from $10^{-7}$ to \SI{e-3}{eV}.

The two-level systems that matter, with asymmetries below a meV and tunneling exponents below about 20, are rare among pairs of neighboring minima. A study of the energy minima of a model glass found of order $10^3$ active two-level systems after sampling some $10^6$ minima~\cite{khomenko2020depletion}. The sampling is classical on the Born--Oppenheimer surface of the potential, $n_B = 0$, with tunneling added through the instanton path. As in that study, glasses of $10^3$ to $10^4$ atoms are quenched with the machine-learned potential, at least $10^4$ of them, with swap Monte Carlo or a comparable method to reach stable glasses. Molecular dynamics at low temperature visits the minima of each glass, pairs of adjacent minima with energy difference below \SI{5}{meV} are kept, and a nudged elastic band connects each pair. The tunneling exponent comes from the instanton path; the deformation potential from the energy difference under a finite strain; the dipole moment from the Berry-phase polarization of the two configurations. To resolve at least $10^3$ systems per decade of splitting, about $10^8$ pairs of minima are screened, some $10^4$ from each glass. The sampling costs $10^5$ to $10^6$ hours on graphics processors; the response tables, under ten core-hours.

The tables contain a principal value and a Hadamard finite part, which are implemented by subtracting the first two terms of the Taylor expansion of the integrand inside a unit interval. The ensemble average is converged to \SI{1}{\percent}, and the finite part is stable to $10^{-3}$ when the inner interval is halved. The measured control is the acoustic loss. The computed tunneling strength $\bar P\gamma^2/\rho_Tc^2$, with $\gamma$ the deformation potential and $c$ the sound speed, must reproduce the measured internal friction of the same glass~\cite{phillips1987two}. The dielectric loss tangent $\tan\delta_0$ of \cref{sec:16.7} gives a second control through the dipole moments. The density of two-level systems in the model glass fell by two orders of magnitude between its least and most stable preparations~\cite{khomenko2020depletion}. A glass quenched on the computer is far less stable than one cooled in the laboratory, so the computed $\bar P$ is meaningful only as a function of the stated preparation. In \cref{sec:16.7,sec:16.8} the tables replace the flat distribution in the scattering and absorption rates, and the distributions of $\Delta_{\mathrm{as}}$, $\Delta_0$ and the couplings test the universality that the standard tunneling model assumes.

\begin{table}[t]
  \caption{Defect and two-level calculations. Costs are estimates. GPU-h are hours on a graphics processing unit (GPU); KKR, Korringa--Kohn--Rostoker; NEB, nudged elastic band; SCF, self-consistent field; MD, molecular dynamics.}
  \label{tab:E22-summary}
  \centering
  \small
  \begin{tabular}{p{0.19\linewidth}p{0.27\linewidth}p{0.29\linewidth}p{0.13\linewidth}}
    \toprule
    quantity & method and grid & acceptance & cost \\
    \midrule
    $\alpha_{\mathrm{ohm}}$ (W, Nb--H, Al, AlO$_x$, WSi) & KKR eigenphases; or supercell overlaps, 54--250 atoms at $\mathrm{\Gamma}$ & Born limit $4K_{\mathrm{b}}$; resonator damping (Al), beyond Born & $10^3$ GPU-h per defect \\
    $S_{\mathrm{HR},k}$, $E_{\mathrm{reorg}}$, $\lambda_{\mathrm{tun}}$ (Si, Ge pairs; Al) & NEB, Hessians, instantons; 64 and 216 atoms; 50 paths & release-path sum rule; $\lambda_{\mathrm{tun}}$ to \SI{5}{\percent} & 120 core-h (pilot); $10^4$--$10^5$ GPU-h \\
    Peierls profile, $M_{\mathrm{cr}}$, $S_{\mathrm{pol}}$ (W); dumbbell glide (Al) & NEB, 250 atoms; MD at 10--\SI{100}{K} & noise $<\SI{0.1}{meV}$ per cell; sinusoidal profile & $10^5$ core-h per material \\
    $\langle|f_{\mathrm{TLS}}|^2\rangle$ (a-SiO$_2$) & two SCFs per system, 216 atoms, \SI{80}{Ry} & \SI{2}{\percent} per system; 100 systems for \SI{10}{\percent} & 200 GPU-h \\
    Two-level ensemble (a-SiO$_2$, AlO$_x$, a-WSi) & $10^8$ pairs of minima; NEB; response tables & \SI{1}{\percent}; acoustic loss & $10^5$--$10^6$ GPU-h \\
    \bottomrule
  \end{tabular}
\end{table}
